\documentclass{article}
\usepackage{tekvideoexercises}

\begin{document}
\exercisename{Kontinuitet}
\streaklength{5}

\tableofcontents
\newpage

\begin{exercise}{Kontinuitet 1-1}

Er funktionen
\[
 f(x) = \begin{cases}
 x + 2, \ x < 1 \\
 3, \ x = 1 \\
 x + 4, \ x > 1
 \end{cases}
\]
kONTINUERT i $x = 1$?

\answerbox{Nej}

\hint

For at en funktion er kontinuert i et punkt, skal grænseværdien fra venstre være lig grænseværdien fra højre og lig funktionsværdien.

\hint

Beregn venstresidet grænse:
\[
\lim_{x \to 1^-} f(x) = \lim_{x \to 1^-} (x + 2) = 1 + 2 = 3
\]

\hint

Beregn højresidet grænse:
\[
\lim_{x \to 1^+} f(x) = \lim_{x \to 1^+} (x + 4) = 1 + 4 = 5
\]

\hint

Vurder funktionsværdien:
\[
f(1) = 3
\]

\hint

Venstresidet grænse = 3, højresidet grænse = 5, funktionsværdi = 3.

\hint

Da venstresidet grænse $\neq$ højresidet grænse, er funktionen ikke kontinuert i $x = 1$.

\end{exercise}

\newpage

\begin{exercise}{Kontinuitet 1-2}

Er funktionen
\[
 g(x) = \begin{cases}
 2x + 1, \ x < 2 \\
 5, \ x = 2 \\
 2x + 1, \ x > 2
 \end{cases}
\]
kONTINUERT i $x = 2$?

\answerbox{Ja}

\hint

Beregn venstresidet grænse:
\[
\lim_{x \to 2^-} g(x) = 2(2) + 1 = 5
\]

\hint

Beregn højresidet grænse:
\[
\lim_{x \to 2^+} g(x) = 2(2) + 1 = 5
\]

\hint

Vurder funktionsværdien:
\[
g(2) = 5
\]

\hint

Alle tre er lig med 5, så funktionen er kontinuert.

\end{exercise}

\newpage

\begin{exercise}{Kontinuitet 1-3}

Er funktionen $h(x) = x^2 + 3x - 4$ kontinuert i $x = 0$?

\answerbox{Ja}

\hint

Polynomier er kontinuerte alle vegne.

\hint

Kontinuitet i $x = 0$ kan bekræftes ved at beregne grænseværdien:
\[
\lim_{x \to 0} h(x) = 0^2 + 3(0) - 4 = -4
\]

\hint

Funktionsværdien i $x = 0$ er: $h(0) = -4$.

\hint

Da grænseværdien eksisterer og er lig funktionsværdien, er funktionen kontinuert.

\end{exercise}

\newpage

\begin{exercise}{Kontinuitet 1-4}

Er funktionen
\[
 p(x) = \frac{x^2 - 1}{x - 1}
\]
kONTINUERT i $x = 1$?

\answerbox{Nej}

\hint

Funktionen er ikke defineret i $x = 1$ (nævneren bliver nul).

\hint

Selvom grænseværdien eksisterer:
\[
\lim_{x \to 1} \frac{x^2 - 1}{x - 1} = \lim_{x \to 1} (x + 1) = 2
\]

\hint

Så er funktionen ikke defineret i $x = 1$ og derfor ikke kontinuert der.

\end{exercise}

\newpage

\begin{exercise}{Kontinuitet 1-5}

Er funktionen $q(x) = |x|$ kontinuert i $x = 0$?

\answerbox{Ja}

\hint

Beregn venstresidet grænse:
\[
\lim_{x \to 0^-} |x| = \lim_{x \to 0^-} (-x) = 0
\]

\hint

Beregn højresidet grænse:
\[
\lim_{x \to 0^+} |x| = \lim_{x \to 0^+} x = 0
\]

\hint

Vurder funktionsværdien:
\[
q(0) = |0| = 0
\]

\hint

Alle tre er lig med 0, så funktionen er kontinuert i $x = 0$.

\end{exercise}

\newpage

\begin{exercise}{Kontinuitet 1-6}

Er funktionen
\[
 r(x) = \begin{cases}
 \frac{1}{x}, \ x \neq 0 \\
 0, \ x = 0
 \end{cases}
\]
kONTINUERT i $x = 0$?

\answerbox{Nej}

\hint

Beregn venstresidet grænse:
\[
\lim_{x \to 0^-} r(x) = \lim_{x \to 0^-} \frac{1}{x} = -\infty
\]

\hint

Beregn højresidet grænse:
\[
\lim_{x \to 0^+} r(x) = \lim_{x \to 0^+} \frac{1}{x} = \infty
\]

\hint

Da grænseværdierne ikke eksisterer (de går mod uendelig), er funktionen ikke kontinuert.

\end{exercise}

\end{document}
